% GENERATED FILE. Edit canonical lessons, then run npm run generate:books.
% canonical-source-hash: fnv1a64:048c6620ba373bd9
% canonical-lessons: RU-C484-gips, RU-C484-kostyl, RU-C484-vyvikh, RU-C484-rastyazheniye, RU-C484-khromat, RU-R484-review-477-484

\chapter{In Plaster}
\label{ch:ru-in-plaster}
\hlchaptermodality{484}

\begin{chapteropening}
\textbf{By the end of this chapter:} \emph{I can talk about broken bones, sprains and limping.}

Complete the last lesson of chapter 484: I can talk about broken bones, sprains and limping.
\end{chapteropening}

\section[\texorpdfstring{\ru{гипс} (gips)}{гипс (gips)}]{\ru{гипс} (gips) — a plaster cast; plaster}

\label{lesson:RU-C484-gips}



\begin{quote}
Before the new one: say the Russian for a graze, then the Russian for to heal.
\end{quote}

\subsection*{You'll want to know: \ru{гипс}}
\textbf{\ru{гипс}} (\emph{gips}) — \textquotedblleft{}a plaster cast; plaster\textquotedblright{}.

\begin{quote}
\ru{Шесть} \ru{недель} \ru{он} \ru{ходил} \ru{с} \ru{рукой} \ru{в} \ru{гипсе}.
\end{quote}

\begin{quote}
\emph{Shest nedél on khodíl s rukóy v gípse.}
\end{quote}

\begin{quote}
He went about with his arm in plaster for six weeks.
\end{quote}

Masculine. \emph{\ru{В} \ru{гипсе}} — in plaster.

\begin{cousinweb}
EtymDB lists \emph{\ru{гипс}} without a source; it is usually traced through German \emph{Gips} to Greek \emph{gýpsos} 'chalk, gypsum'.
\end{cousinweb}

\begin{grammarlens}[title={What you've built}]
One more word for telling, explaining and describing at B1.
\end{grammarlens}

\subsection*{Your turn}
\begin{itemize}
\raggedright
  \item \emph{Say it:} \emph{gips}
  \item \emph{Say it:} \emph{gips}, once more
  \item \emph{Say it:} a sentence of your own with \emph{gips}
  \item \emph{Recall:} say the Russian for a graze, then the Russian for to heal, then say \emph{gips} again
\end{itemize}

\begin{tcolorbox}[breakable,colback=teal!4,colframe=teal!35!black,title={Before you move on}]
What does \ru{гипс} mean? (\textquotedblleft{}A plaster cast; plaster\textquotedblright{}.) Say it once more.
\end{tcolorbox}
\section[\texorpdfstring{\ru{костыль} (kostýl)}{костыль (kostýl)}]{\ru{костыль} (kostýl) — a crutch}

\label{lesson:RU-C484-kostyl}



\begin{quote}
Before the new one: say the Russian for to heal, then the Russian for a plaster cast.
\end{quote}

\subsection*{You'll want to know: \ru{костыль}}
\textbf{\ru{костыль}} (\emph{kostýl}) — \textquotedblleft{}a crutch\textquotedblright{}.

\begin{quote}
\ru{После} \ru{операции} \ru{бабушка} \ru{ходила} \ru{на} \ru{костылях}.
\end{quote}

\begin{quote}
\emph{Pósle operátsii bábushka khodíla na kostylyákh.}
\end{quote}

\begin{quote}
After the operation Grandma walked on crutches.
\end{quote}

Masculine; genitive \ru{костыля}. \emph{\ru{Ходить} \ru{на} \ru{костылях}} — to walk on crutches.

\begin{cousinweb}
\emph{\ru{костыль}} is derived from \emph{\ru{кость}} 'bone' (EtymDB), which is inherited from Proto-Slavic \emph{kost\ru{ь}}.
\end{cousinweb}

\begin{grammarlens}[title={What you've built}]
One more word for telling, explaining and describing at B1.
\end{grammarlens}

\subsection*{Your turn}
\begin{itemize}
\raggedright
  \item \emph{Say it:} \emph{kostýl}
  \item \emph{Say it:} \emph{kostýl}, once more
  \item \emph{Say it:} a sentence of your own with \emph{kostýl}
  \item \emph{Recall:} say the Russian for to heal, then the Russian for a plaster cast, then say \emph{kostýl} again
\end{itemize}

\begin{tcolorbox}[breakable,colback=teal!4,colframe=teal!35!black,title={Before you move on}]
What does \ru{костыль} mean? (\textquotedblleft{}A crutch\textquotedblright{}.) Say it once more.
\end{tcolorbox}
\section[\texorpdfstring{\ru{вывих} (vývikh)}{вывих (vývikh)}]{\ru{вывих} (vývikh) — a dislocation}

\label{lesson:RU-C484-vyvikh}



\begin{quote}
Before the new one: say the Russian for a plaster cast, then the Russian for a crutch.
\end{quote}

\subsection*{You'll want to know: \ru{вывих}}
\textbf{\ru{вывих}} (\emph{vývikh}) — \textquotedblleft{}a dislocation\textquotedblright{}.

\begin{quote}
\ru{Врач} \ru{сказал}, \ru{что} \ru{у} \ru{меня} \ru{вывих} \ru{плеча}.
\end{quote}

\begin{quote}
\emph{Vrach skazál, chto u menyá vývikh plechá.}
\end{quote}

\begin{quote}
The doctor said I had a dislocated shoulder.
\end{quote}

Masculine. The verb is \emph{\ru{вывихнуть}} — to dislocate, to twist.

\begin{cousinweb}
\emph{\ru{вывих}} is built on \emph{\ru{вы}-} 'out' and the stem of \emph{\ru{вывихнуть}} 'to dislocate'; EtymDB has no entry, and the origin of the stem is not settled.
\end{cousinweb}

\begin{grammarlens}[title={What you've built}]
One more word for telling, explaining and describing at B1.
\end{grammarlens}

\subsection*{Your turn}
\begin{itemize}
\raggedright
  \item \emph{Say it:} \emph{vývikh}
  \item \emph{Say it:} \emph{vývikh}, once more
  \item \emph{Say it:} a sentence of your own with \emph{vývikh}
  \item \emph{Recall:} say the Russian for a plaster cast, then the Russian for a crutch, then say \emph{vývikh} again
\end{itemize}

\begin{tcolorbox}[breakable,colback=teal!4,colframe=teal!35!black,title={Before you move on}]
What does \ru{вывих} mean? (\textquotedblleft{}A dislocation\textquotedblright{}.) Say it once more.
\end{tcolorbox}
\section[\texorpdfstring{\ru{растяжение} (rastyazhéniye)}{растяжение (rastyazhéniye)}]{\ru{растяжение} (rastyazhéniye) — a strain (of a muscle)}

\label{lesson:RU-C484-rastyazheniye}



\begin{quote}
Before the new one: say the Russian for a crutch, then the Russian for a dislocation.
\end{quote}

\subsection*{You'll want to know: \ru{растяжение}}
\textbf{\ru{растяжение}} (\emph{rastyazhéniye}) — \textquotedblleft{}a strain (of a muscle)\textquotedblright{}.

\begin{quote}
\ru{У} \ru{футболиста} \ru{растяжение} \ru{мышцы}.
\end{quote}

\begin{quote}
\emph{U futbolísta rastyazhéniye mýshtsy.}
\end{quote}

\begin{quote}
The footballer has a pulled muscle.
\end{quote}

Neuter. \emph{\ru{Растяжение} \ru{связок}} — a sprain.

\begin{cousinweb}
\emph{\ru{растяжение}} is built on \emph{\ru{растянуть}} 'to stretch', from \emph{\ru{тянуть}} 'to pull', which is inherited from Proto-Slavic \emph{tęgnǫti} (EtymDB).
\end{cousinweb}

\begin{grammarlens}[title={What you've built}]
One more word for telling, explaining and describing at B1.
\end{grammarlens}

\subsection*{Your turn}
\begin{itemize}
\raggedright
  \item \emph{Say it:} \emph{rastyazhéniye}
  \item \emph{Say it:} \emph{rastyazhéniye}, once more
  \item \emph{Say it:} a sentence of your own with \emph{rastyazhéniye}
  \item \emph{Recall:} say the Russian for a crutch, then the Russian for a dislocation, then say \emph{rastyazhéniye} again
\end{itemize}

\begin{tcolorbox}[breakable,colback=teal!4,colframe=teal!35!black,title={Before you move on}]
What does \ru{растяжение} mean? (\textquotedblleft{}A strain (of a muscle)\textquotedblright{}.) Say it once more.
\end{tcolorbox}
\section[\texorpdfstring{\ru{хромать} (khromát)}{хромать (khromát)}]{\ru{хромать} (khromát) — to limp}

\label{lesson:RU-C484-khromat}



\begin{quote}
Before the new one: say the Russian for a dislocation, then the Russian for a strain.
\end{quote}

\subsection*{You'll want to know: \ru{хромать}}
\textbf{\ru{хромать}} (\emph{khromát}) — \textquotedblleft{}to limp\textquotedblright{}.

\begin{quote}
\ru{Собака} \ru{хромает} \ru{на} \ru{переднюю} \ru{лапу}.
\end{quote}

\begin{quote}
\emph{Sobáka khromáyet na perédnyuyu lápu.}
\end{quote}

\begin{quote}
The dog is limping on its front paw.
\end{quote}

Imperfective; present \ru{хромаю}. The adjective is \emph{\ru{хромой}} — lame.

\begin{cousinweb}
\emph{\ru{хромать}} is built on \emph{\ru{хромой}} 'lame'; EtymDB lists both without a source, and \emph{\ru{хромой}} is usually traced to Proto-Slavic \emph{xrom\ru{ъ}}.
\end{cousinweb}

\begin{grammarlens}[title={What you've built}]
One more word for telling, explaining and describing at B1.
\end{grammarlens}

\subsection*{Your turn}
\begin{itemize}
\raggedright
  \item \emph{Say it:} \emph{khromát}
  \item \emph{Say it:} \emph{khromát}, once more
  \item \emph{Say it:} a sentence of your own with \emph{khromát}
  \item \emph{Recall:} say the Russian for a dislocation, then the Russian for a strain, then say \emph{khromát} again
\end{itemize}

\begin{tcolorbox}[breakable,colback=teal!4,colframe=teal!35!black,title={Before you move on}]
What does \ru{хромать} mean? (\textquotedblleft{}To limp\textquotedblright{}.) Say it once more.
\end{tcolorbox}
\section[(dialogue)]{Words from chapters 477–484 — a review}

\label{lesson:RU-R484-review-477-484}



\begin{quote}
Say the Russian for longing and to pine.
\end{quote}

\subsection*{Your turn}
Say these aloud:

\begin{itemize}
\raggedright
  \item longing, to pine, loneliness, lonely, to regret
  \item stubborn, stubbornness, fussy, to play up, to grumble
  \item trust, to trust, to betray, to respect, sincere
  \item to roar with laughter, to giggle, to sob, to comfort, consolation
  \item a shoulder blade, the spine, a thigh, a heel, a palm
  \item a temple, the back of the head, an eyelid, a nostril, a jaw
  \item a scratch, to hurt oneself, a cut, a graze, to heal
  \item a plaster cast, a crutch, a dislocation, a strain, to limp
\end{itemize}

\begin{tcolorbox}[breakable,colback=teal!4,colframe=teal!35!black,title={Before you move on}]
Longing? (\textbf{\ru{тоска}}.) To pine? (\textbf{\ru{тосковать}}.) Loneliness? (\textbf{\ru{одиночество}}.) Lonely? (\textbf{\ru{одинокий}}.) To regret? (\textbf{\ru{сожалеть}}.) Stubborn? (\textbf{\ru{упрямый}}.) Stubbornness? (\textbf{\ru{упрямство}}.) Fussy? (\textbf{\ru{капризный}}.) To play up? (\textbf{\ru{капризничать}}.) To grumble? (\textbf{\ru{ворчать}}.) Trust? (\textbf{\ru{доверие}}.) To trust? (\textbf{\ru{доверять}}.) To betray? (\textbf{\ru{предавать}}.) To respect? (\textbf{\ru{уважать}}.) Sincere? (\textbf{\ru{искренний}}.) To roar with laughter? (\textbf{\ru{хохотать}}.) To giggle? (\textbf{\ru{хихикать}}.) To sob? (\textbf{\ru{рыдать}}.) To comfort? (\textbf{\ru{утешать}}.) Consolation? (\textbf{\ru{утешение}}.) A shoulder blade? (\textbf{\ru{лопатка}}.) The spine? (\textbf{\ru{позвоночник}}.) A thigh? (\textbf{\ru{бедро}}.) A heel? (\textbf{\ru{пятка}}.) A palm? (\textbf{\ru{ладонь}}.) A temple? (\textbf{\ru{висок}}.) The back of the head? (\textbf{\ru{затылок}}.) An eyelid? (\textbf{\ru{веко}}.) A nostril? (\textbf{\ru{ноздря}}.) A jaw? (\textbf{\ru{челюсть}}.) A scratch? (\textbf{\ru{царапина}}.) To hurt oneself? (\textbf{\ru{ушибиться}}.) A cut? (\textbf{\ru{порез}}.) A graze? (\textbf{\ru{ссадина}}.) To heal? (\textbf{\ru{заживать}}.) A plaster cast? (\textbf{\ru{гипс}}.) A crutch? (\textbf{\ru{костыль}}.) A dislocation? (\textbf{\ru{вывих}}.) A strain? (\textbf{\ru{растяжение}}.) To limp? (\textbf{\ru{хромать}}.)
\end{tcolorbox}
